\documentclass[conference,a4paper]{IEEEtran}
\usepackage{cite}
\usepackage{amsmath,amssymb,amsfonts}
\usepackage{graphicx}
\usepackage{textcomp}
\usepackage{xcolor}
\usepackage{hyperref}
\usepackage{lettrine}

\bibliographystyle{unsrt}

\IEEEoverridecommandlockouts
\begin{document}

\title{Backside Power Delivery for Sub-5 nm CMOS: Signal Integrity Benefits and Emerging Design Trade-offs}

\author{\IEEEauthorblockN{Daniel Tyukov (5714699)\thanks{This article was written as part of the course EE4C01 Profile Orientation \& Academic Skills of the faculty Electrical Engineering, Mathematics and Computer Science of Delft University of Technology. Word count: 1535.}}
\IEEEauthorblockA{Department of Microelectronics\\
Delft University of Technology\\
Delft, The Netherlands\\
d.tyukov@student.tudelft.nl}
\and
\IEEEauthorblockN{Apostolos Tsailas (5506387)}
\IEEEauthorblockA{Department of Microelectronics\\
Delft University of Technology\\
Delft, The Netherlands\\
a.tsailas-1@student.tudelft.nl}
}

\maketitle

\begin{abstract}
    Technology scaling beyond 5nm for CMOS has revealed major limitations in conventional front-side power delivery networks (FSPDNs), including increased interconnect resistance and routing congestion. In the FSPDN architecture, power and signal interconnects share the same resources, making it increasingly difficult to maintain power integrity without impacting performance. This has put the spotlight on a different architecture, back-side power delivery networks (BSPDNs). 
    
    This paper review first outlines key limitations of FSPDNs and then examines the benefits of relocating the power network to the backside of the wafer. BSPDNs decouple signal and power interconnect, which enables a reduction in peak IR drop of up to 85\%, a cell area reduction of around 20\% and other performance improvements.
    
    The review also highlights challenges associated with the BSPDN architecture, including increased thermal resistance and increased temperature in hotspots. Furthermore, an increased fabrication complexity due to accompanying features such as nano-TSVs and buried power rails (BPR). 

    Overall, BSPDNs have demonstrated to be a promising evolution in power delivery architecture which will enable scaling beyond the 5nm range.
\end{abstract}

\section{Introduction}

\lettrine{T}{echnology} scaling beyond 5nm has introduced significant challenges in power delivery and interconnect performance. Although transistor dimensions continue to shrink, the scaling of interconnects has not kept pace, leading to increased wire resistance and tighter voltage margins. Conventional front-side power delivery networks (PDNs), which share routing resources with signal interconnects, face growing difficulty in meeting power integrity requirements without degrading routing efficiency and overall performance\cite{8993617}.

To address such limitations buried power rails (BPR) have been proposed as an early innovation. By relocating power rails beneath active devices and utilizing high-aspect-ratio ruthenium structures, BPR reduces local resistance and frees front-side routing tracks, enabling improved IR-drop characteristics\cite{8993617}. Backside Power Delivery Networks (BS-PDN) move the power grid to the wafer backside which decouples the power and signal routing\cite{10185208}. Experimental findings have shown substantial reductions in voltage droop and measurable frequency improvements in high-performance cells \cite{10185208}. It was also demonstrated significant area savings and enhanced routing efficiency were achieved in both high-performance and high-density scenarios\cite{10631463}, while models indicate more than a four times reduction in IR-drop compared to conventional PDN\cite{8932458}.

These architectural innovations represent a fundamental shift in power delivery design, moving from incremental BEOL optimizations toward structural reorganization of the power distribution network. This literature review therefore examines the evolution from conventional front-side PDNs to buried power rails and backside power delivery, evaluating their signal integrity and power–performance–area benefits as well as the emerging design challenges they introduce in sub-5nm CMOS technologies.
\section{Methodology}

Literature was collected from IEEE Xplore. Search terms combined ``backside power delivery,'' ``buried power rails,'' ``front-side PDN,'' ``IR-drop,'' ``signal integrity,'' and ``sub-5nm CMOS.'' More specific terms such as ``nanoTSV,'' ``PowerVia,'' and ``power-performance-area'' were added to capture implementation studies. Only publications from 2019 onward were considered, since the earliest demonstrations of buried power rails and backside delivery at advanced nodes date from that year.

The initial search returned over 80 results. Titles and abstracts were screened first: papers had to deal directly with power delivery network architecture and its effect on power integrity or routing efficiency at sub-5\,nm nodes. Work focused purely on transistor-level scaling or on unrelated packaging topics was excluded. The remaining papers were then read in full. Each was judged on whether its results rested on simulation, measurement, or both, and on how much it added to the comparison between front-side and backside delivery. Papers that reported quantitative IR-drop, voltage droop, or area figures across different PDN configurations were preferred.

Nine sources were retained. They span early architecture proposals, industry demonstrations from Intel and imec, and parasitic modelling of backside configurations. Together they cover the path from conventional front-side PDN through buried power rails to full backside delivery, with both measured and simulated data on the trade-offs involved.

\section{Limitations of Conventional Front-Side Power Delivery Networks}

Conventional integrated circuits rely on front-side power delivery (FSPDNs) where power and signal interconnects are routed via the same back-end-of-line (BEOL) metal stack. In this layout, supply voltages are distributed from the package connections through all individual metal layers, down to the transistor level. This architecture has been used for multiple generations of integrated circuits, however continued scaling below 5 nm has revealed several limitations that have made reliable power delivery more difficult \cite{9817394},\cite{10244504}. 

As a result of the scaling of interconnects, an issue arises from the increased resistance. As metal dimensions shrink, resistivity increases due to surface and grain-boundary scattering, leading to larger "IR" drops across the power grid\cite{8993617}. This results in a reduced effective supply voltage at the transistors, which can degrade circuit performance and reliability. In the most advanced node, highly resistive local metal layers contribute significantly to IR-drop hotspots, particularly in areas with dense logic \cite{8993617}. Studies conducted with an ARM Cortex-A53 design equivalent to a 3nm technology node show that conventional FSPDNs did not meet the IR-drop targets, highlighting the limitations of this architecture at sub-5 nm technology node \cite{9817394}.

Another key limitation that arises due to scaling is the increased competition for the limited routing resources available for power and signal interconnects. In FSPDNs, both power rails and signal interconnects occupy the same metal stack. As such, designers must dedicate a large portion of routing resources to maintain a sufficiently reliable power grid that is able to provide the necessary performance. This results in reduced routing resources for signal interconnects and increased routing congestion. Experiments have shown that high-performance designs often require dense FSPDNs to achieve the necessary IR drop target, which consumes valuable routing tracks, resulting in designs that occupy a larger area \cite{10631463}.

To mitigate these limitations, Buried Power Rails (BPR) are used as a current solution. These designs move power rails beneath active devices in order to reduce resistance and make more routing resources available to the signal interconnect. BPR-based designs have shown a 25\% reduction in on-chip IR drop and 17\% reduction in off-chip voltage droop compared with conventional front-side PDNs \cite{9817394}. However, even with these improvements, standalone BPR designs are viewed as intermediate solutions\cite{10244504}. The fundamental limitations of front-side routing remain.

Backside power delivery networks (BSPDNs) have been explored as a long-term solution to these issues. In this architecture, the power grid is moved to the backside of the wafer. By separating power and signal routing, BSPDN architecture reduces routing congestion and allows for reliable power transmission. This increases the scalability of future technology nodes\cite{10631463},\cite{8932458}.


\section{Improvements of Back-Side Power Delivery Network Implementations}

Back-Side Power Delivery Networks improve power integrity and routing efficiency by relocating the power grid to the backside of the wafer and utilising the front-side solely for signal interconnect. The design allows for a reduction in congestion and increases routing flexibility, especially in advanced technology nodes\cite{10631463}. In dense logic areas, this architecture allows for signal interconnects to utilize lower layers on the front side of the wafer and, as a result, reduces wire lengths, reduces interconnect delay, and improves timing\cite{11192533}.

A primary improvement of the BSPDN architecture is the reduction in IR drop. By delivering power from a dedicated backside metal layer, voltage droop at the transistor level is reduced. Modelling of backside PDN configurations incorporating buried power rails and µTSVs demonstrated a more than four-fold reduction in IR-drop compared to conventional front-side PDNs\cite{8932458}. This improvement is due to the lower resistance of the thicker backside power rails and the shorter current path to the transistor layer, which bypasses the highly resistive local BEOL metal layers entirely.

Intel's PowerVia Technology, demonstrated using a fabricated E-core, showed that a BSPDN defforsign yielded a 30\% improvement in voltage droop and 6\% frequency improvement compared to a reference FSPDN design\cite{10185208}. By splitting power delivery from the signal interconnect, PowerVia supplies a more uniform supply voltage across the die, resulting in lower average and worst-case IR-drop hotspots that are particularly problematic in dense logic regions\cite{10185208}.

BSPDN architecture has also shown improvements in routing efficiency and cell area. With the power grid moved to the backside of the wafer, front-side routing tracks that would be used for power delivery are freed, enabling either area reduction or increased routing flexibility at the same cell dimensions. Studies evaluating high-density designs showed area savings up to 19\% when implementing a BSPDN architecture. Signal interconnect was completed at lower metal layers, reducing overall wire length\cite{10631463}.

Table~\ref{tab:comparison} summarises the quantitative improvements reported across the reviewed studies for each power delivery architecture relative to the conventional FSPDN baseline.

\begin{table}[h]
\centering
\caption{Comparison of power delivery architectures at sub-5\,nm nodes. Values indicate improvement relative to conventional FSPDN.}
\label{tab:comparison}
\begin{tabular}{|l|c|c|c|}
\hline
\textbf{Metric} & \textbf{FSPDN} & \textbf{BPR} & \textbf{BSPDN} \\
\hline
On-chip IR-drop & Baseline & $-$25\%\cite{9817394} & $>$4$\times$ reduction\cite{8932458} \\
\hline
Off-chip voltage droop & Baseline & $-$17\%\cite{9817394} & $-$30\%\cite{10185208} \\
\hline
Frequency & Baseline & --- & +6\%\cite{10185208} \\
\hline
Cell area & Baseline & --- & $-$19\%\cite{10631463} \\
\hline
Thermal (hotspot) & Baseline & --- & +45\%\cite{10925422} \\
\hline
\end{tabular}
\end{table}
\section{Limitations of Back-Side Power Delivery Networks}

The improvements in power integrity and routing efficiency offered by BSPDNs come at a cost. Moving the power grid to the backside of the wafer introduces thermal, self-heating, and AC impedance challenges that are not present in conventional front-side architectures.

The most pressing issue is thermal performance. In FSPDNs, the full-thickness silicon substrate spreads heat laterally from the active devices toward the topside heat sink. BSPDNs require thinning the wafer to form nano through-silicon via (nTSV) connections, which removes much of this lateral heat-spreading path. The bonding interface and back-end-of-line (BEOL) layers between the devices and the heat sink further increase thermal resistance. Modelling of CPU hotspot regions has shown that BSPDN configurations reach temperatures approximately 45\% higher than equivalent FSPDN designs\cite{10925422}. Several mitigation strategies have been investigated, including high thermal conductivity bonding materials, stacked BEOL via configurations that create direct vertical thermal paths, and embedded microchannel cooling etched into the backside metal layers\cite{10925422}.

At the device level, the thinned substrate also worsens self-heating effects (SHE). With less silicon to absorb and spread the heat generated by switching transistors, localised temperature rises become more severe. An evaluation of SHE in SRAM macros and circuits at the 3nm node found that BSPDN designs improve power, performance, and area, but the resulting thermal behaviour requires power-thermal co-design to avoid reliability degradation\cite{10529350}. Electrical gains therefore cannot be evaluated in isolation from their thermal consequences.

A separate concern relates to the AC impedance profile of the backside power grid. Because the DC resistance of BSPDNs is much lower than that of FSPDNs, impedance resonance peaks can shift into frequency ranges that coincide with circuit switching activity, risking supply voltage fluctuations\cite{8993617}. Addressing this requires decoupling capacitor placement and impedance profiling during design.


\section{Conclusion}
This literature review has examined the limitations of conventional front-side PDNs, the advantages of back-side PDNs, and the challenges introduced by these architectures in advanced technology nodes. Scaling beyond the 5nm range, increased interconnect resistance,  reduced voltage margins and routing congestion have made traditional front-side PDNs increasingly inefficient and difficult to optimise. 

The reviewed studies demonstrate that back-side PDNs provide substantial improvements in power integrity and design efficiency. By relocating the power network to the backside of the wafer, this architecture decouples power and signal lines, thereby increasing routing efficiency and better utilizing front-end metal layers. Major reductions in IR drop have been demonstrated, with reported improvements of up to 85\% compared to front-side PDNs designs. Furthermore, back-side implementations have shown clear gains in power, performance, and area, including frequency improvements of 6\% and core area reductions of about 20\% and upwards depending on the design and technology scenario.

Despite these advantages, back-side PDNs introduce new challenges. The relocation of the power network affects heat dissipation, resulting in higher thermal resistance and hotspots with a 45\% increase in temperature. Furthermore, integration of additional architectural features such as Buried Power Rails and nano-TSVs introduces additional parasitic effects that need to be carefully managed.

In conclusion, back-side PDNs demonstrate a significant evolution in power delivery architecture that addresses limitations of conventional front-side PDNs. Although back-side PDNs introduce new challenges in thermal management and fabrication, the architecture's improvements in IR drop, performance and cell area make it a strong candidate for technology nodes in the future. While further research in these limitations is necessitated, current evidence indicated that back-side PDNs will have a significant role in sub-5nm CMOS design.

%TC:ignore 
\bibliography{bibliography}
 %TC:endignore

\end{document}
